\section{The research question and the scope of the answer}
\label{sec:scope}

The inspected ProveIt manuscript develops exact identities among
polylogarithms, Gamma derivatives, harmonic tails, and Hurwitz--Stieltjes
coordinates \cite{RepoMain}. Its differential-polynomial continuation
already proves a canonical integration-by-parts normal form and a complete
recurrence for the variables $u_0,u_1,u_2$ \cite{RepoQuadratic}. A later
continuation treats square sectors on an affine background and explicitly
asks for every coefficient
\begin{equation}\label{eq:source-question}
 [u_0^a u_1^k u_2^b u_3^2]N_{a+2k+3b+8},\qquad b>0,
\end{equation}
with parity and the endpoint primitive retained \cite[Question 1]{RepoQuestions}.
Here $u_r$ represents the $r$th derivative of digamma, and $N_m$ is the
canonical normal representative of
$(-1)^m\Gamma(x)(\mathrm d/\mathrm dx)^m\Gamma(x)^{-1}$.

We answer \eqref{eq:source-question} without restricting $a,k,b$. More
strongly, we allow any number $n\ge2$ of $u_3$ factors. The answer is a
finite rational sum whose summands have four consecutive odd factors in
the denominator. Its underlying integral is a beta moment with the
endpoint factor $(1-v^2)^3$. This factor is not decorative: its three
vanishing endpoint derivatives are exactly what makes the third-order
variational transport close.

\subsection{A preview in ordinary coefficients}
For $m=a+2k+3b+4n$ and $n\ge2$, the coefficient vanishes when $k+n$ is
odd. When $k+n$ is even it is
\begin{equation}\label{eq:preview}
\begin{split}
 &[u_0^a u_1^k u_2^b u_3^n]N_m\\
 &\quad=-\frac{m!}{192\,a!\,2^k k!\,b!\,n!\,24^b48^{n-2}}
 \sum_{j=0}^{b}\sum_{\ell=0}^{n-2}
 \frac{\binom bj3^j\binom{n-2}{\ell}}
 {\displaystyle\prod_{r=0}^{3}(k+n+2j+2\ell-1+2r)}.
\end{split}
\end{equation}
For example, the coefficients of $u_2u_3^2$, $u_2^2u_3^2$, and
$u_1u_3^3$ are respectively $-55$, $-3640$, and $-3185/3$. These are
coefficients in the specified rational normal basis, not assertions of
arithmetic independence after special-function substitution.

\subsection{What is inherited, and what is established here}
The normal-form decomposition and the finite-jet Euler principle are
inherited. We reproduce their elementary proofs so that the article does
not depend on trusting a computation or an inaccessible manuscript.
The underlying functional-monomial framework is credited in the source to
Guo, Regensburger, and Rosenkranz \cite{GRR}. The preceding pure-$u_2$
recurrence is also inherited; its beta-density solution and the complete
cubic-background formula are the developments proved below.

The polynomial-weight Hurwitz and Lerch integration-by-parts mechanism
uses classical special-function identities \cite{DLMF25,DLMF14}. We do not
claim that mechanism, the standard Gamma--zeta relations, or elementary
partial fractions as new. Their application to the newly determined mixed
coefficient kernels gives the exact families and concrete evaluations in
Sections~\ref{sec:hurwitz} and \ref{sec:polylog}. Historical priority for
these specializations has not been exhaustively investigated.

The incoming directory was inventoried during source selection, but its
binary packages were not exhaustively read or replayed. In particular,
the conclusions here are not a claim that every current delivery has been
reconciled. The provenance file distinguishes the inspected text from
archive names seen in the directory. Since the repository changed during
retrieval, we record file-level identifiers where available instead of
claiming a single atomic repository snapshot.

\subsection{Conventions}
All formal polynomial identities are over $\Q$. We write
\[
 \avg{f(v)}=\frac12\int_{-1}^1f(v)\dd v,
 \qquad [t^m]F(t)\text{ for ordinary coefficient extraction}.
\]
The exponential generating functions carry their displayed factorials.
The notation $\zeta^{(r)}(s,a)$ always means an \emph{order} derivative,
$\partial_s^r\zeta(s,a)$; $\psi^{(r)}(a)$ always means an \emph{argument}
derivative. Stieltjes constants use the convention
\begin{equation}\label{eq:stieltjes-convention}
 \zeta(1+\epsilon,a)=\frac1\epsilon+
 \sum_{q\ge0}\frac{(-1)^q}{q!}\gamma_q(a)\epsilon^q,
 \qquad \gamma_0(a)=-\psi(a).
\end{equation}
For analytic integral formulas we initially impose $\Re a>1$, so that
$a+v$ stays in the right half-plane. All polylogarithm series are initially
used on $|z|<1$, eliminating branch ambiguity. Extensions and removable
singularities are stated where used.
